\begin{abstract}
Large Language Model (LLM) agents are increasingly entrusted with autonomous tool invocation across enterprise, financial, and civic administration systems. However, recent empirical security benchmarks reveal that agents frequently exhibit \textit{over-privileged tool selection}---invoking high-permission APIs (e.g., executing raw database queries or global exports) when lower-permission alternatives (e.g., single-record lookups) strictly suffice to fulfill user intent. While this vulnerability has been studied in English, the security implications of \textit{cross-lingual transfer} into regional Southeast Asian languages remain unexplored. In this paper, we investigate the susceptibility of open-weight multilingual LLMs (including SEA-LION, Qwen~2.5, and Llama~3.1) and frontier models (GPT-4o) to over-privileged tool selection when prompted in English, Indonesian, and Javanese (\textit{Basa Jawa}). We demonstrate that morphological indirectness and intent ambiguity in regional languages significantly exacerbate over-privilege rates ($OPR$, surging up to $58.3\%$ in Javanese). To mitigate this vulnerability, we propose a \textit{deterministic, lattice-based least-privilege guardrail} that intercepts candidate tool emissions at runtime, deterministically inspects parameter boundaries, and automatically down-casts excessive privileges into minimal-sufficient equivalents while preserving functional capabilities. Across 432 controlled experimental executions, our guardrail eliminates over-privileged invocations ($OPR = 0.0\%$), blocks transient panic escalation ($PER = 0.0\%$), and preserves perfect task completion utility ($TCR = 100.0\%$) with sub-15ms execution overhead.
\end{abstract}
